\subsection{The unramified obstruction in every weight and the class-four image}

\begin{statement}{A.204}{Proved}
{General strict residue containment and a locally matching class-three lifting theorem}
For \(r\geq2\), put \(E_r=E_U/\Gamma_{r+1}\Pi\), with
closed lower-central subgroups.  Let \(s\) be an actual continuous
section of \(E_r\to G_\Q\) which, at every finite place, is
geometrically conjugate to the image of a full point or cusp
section.  Let \(b_{r+1}\) be the number of weight-\(r+1\)
Hall commutators on two generators.  For every coefficient
prime \(\ell\), its actual next-layer obstruction satisfies
\[
\begin{split}
 o_{r+1,\ell}(s)\ \in\
 \operatorname{im}\bigl(&H^2_{\rm cts}
       (G_{\Q,\{\ell,\infty\}},\Z_\ell(r+1)^{b_{r+1}})\\
       &\longrightarrow H^2_{\rm cts}
       (G_\Q,\Z_\ell(r+1)^{b_{r+1}})\bigr).
\end{split}
\tag{A.204.1}
\]
In particular the complete full-coefficient obstruction is
annihilated by the explicit finite integer
\[
                 B_{r+1}=2\,|K_{2r}(\Z)|.
\tag{A.204.2}
\]
The strict residues have actual primitives tending to zero;
thus (A.204.1) includes the coefficient-power derived-limit
part and uses no finite-support hypothesis on the valuation
coordinates.

For \(r=3\), assume also that the real restriction is the
image of a full real point or cusp section.  Then
\[
 \mathscr S^{\rm loc,all}_3\ \subseteq\
 \operatorname{im}\bigl(\operatorname{Sec}(E_4/G_\Q)
          \longrightarrow\operatorname{Sec}(E_3/G_\Q)\bigr),
\tag{A.204.3}
\]
where the domain on the left consists precisely of the
specified finite and real local comparisons.
This gives a global continuous class-four lift of each such
class-three section.  It does not assert that the new lift
has the specified local point or cusp images.
\end{statement}

\begin{proof}
\emph{The actual central extension and set lift.}
Retain the rational tangential splitting and generators
\(x,y\) of A.201.  The integral Hall construction of
(A.65.2) gives actual continuous geometric coordinates,
and its weight-\(r+1\) kernel is
\[
 Z_{r+1,\ell}=\Gamma_{r+1}\Pi_\ell/\Gamma_{r+2}\Pi_\ell
       =\Z_\ell(r+1)^{b_{r+1}}.
\tag{A.204.4}
\]
The primitive Magnus coefficient matrices there prove this
identification integrally.  The geometric action on this
central kernel is trivial; the arithmetic action is
\(\chi_\ell^{r+1}\), obtained by bracketing the actual
degree-one action.  Write \(N_{r,\ell}\) for the pro-\(\ell\)
geometric class-\(r\) quotient.  The continuous map
\[
 L_r:N_{r,\ell}\longrightarrow N_{r+1,\ell}
\tag{A.204.5}
\]
keeps every Hall coordinate of weight at most \(r\) and
sets the weight-\(r+1\) coordinates to zero.  It is an
actual set section of the projection \(p\).
For the geometric cocycle \(n(g)\) of \(s\), set
\[
 R(g,h)=L_r(n(g))\,gL_r(n(h))\,L_r(n(gh))^{-1}.
\tag{A.204.6}
\]
It is a continuous normalized two-cocycle with values in
(A.204.4), representing exactly the stated lifting obstruction.

\emph{Uniform continuity on all integral conjugators.}
At \(q\ne\ell\) the pointed action on every geometric
pro-\(\ell\) quotient is unramified.  The smooth-pair and
unit-tangent comparison used in A.201 proves this on every
finite \(\ell\)-quotient; alternatively the exact induction
through the automorphism kernels in A.65 proves it at every
Hall weight.  The formal-disc maps of A.49 and A.201 then
give the same three inertia elements
\(c=x,y,(xy)^{-1}\) in all these quotients.

Align an actual local full point or cusp section with the
prescribed \(s\) in \(E_{r,\ell}\).  Lifting its geometric
conjugator through \(N_{r+1,\ell}\to N_{r,\ell}\) gives a
local section \(\sigma_q\) whose projection equals
\(s|_{G_{\Q_q}}\) literally.  In normalized tame inertia
its geometric value is
\[
          u_qc_q^{m_q}u_q^{-1},\qquad
 u_q\in N_{r+1,\ell},\quad m_q\in\Z_\ell,
\tag{A.204.7}
\]
or the identity.  Its abelian coordinates are the same
four patterns as (A.201.7).  Reducing the two actual
global abelian cocycles modulo \(\ell^a\) therefore
proves \(m_q\to0\) as \(q\) leaves finite sets.

For each of the three fixed \(c\), define the exact map
\[
 \mathcal F_{r,c}(u,m)=
 u c^m u^{-1}L_r\bigl(p(u c^m u^{-1})\bigr)^{-1}
                  \in Z_{r+1,\ell}.
\tag{A.204.8}
\]
It is continuous, and \(\mathcal F_{r,c}(u,0)=0\) for
every \(u\).  Its uniform consequence follows from the
actual compact domain, without any assumption on the
chosen local conjugators.  For each \(a\), the inverse
image of \(\ell^aZ_{r+1,\ell}\) is an open neighborhood
of \(N_{r+1,\ell}\times\{0\}\).  A finite subcover of
the compact first factor gives an integer \(d\) such that
\[
 m\in\ell^d\Z_\ell\quad\Longrightarrow\quad
 \mathcal F_{r,c}(u,m)\in\ell^aZ_{r+1,\ell}
                    \quad\hbox{for every }u.
\tag{A.204.9}
\]
Taking the maximum of the three integers covers all cusp
labels.  Thus \(\mathcal F_{r,c_q}(u_q,m_q)\to0\).
This is a statement about the literal Hall map; it does
not replace a cohomology group by a separated quotient.

Write the aligned local section uniquely as
\(\sigma_q(g)=(z_q(g)L_r(n(g)),g)\).
Its homomorphism identity gives \(\mathrm dz_q=-R\),
and evaluation at the same tame inertia generator gives
\[
                  z_q(T_q)=\mathcal F_{r,c_q}(u_q,m_q).
\tag{A.204.10}
\]
For an integral ordinary point it is zero.  It follows
from (A.204.9) that \((z_q(T_q))_{q\ne\ell}\) is an
actual family tending to zero.

\emph{The strict residue and the unramified image.}
For weight \(r+1\), use the same strict cochain map as
(A.201.14), with the same arithmetic Frobenius and tame
generator at every coefficient level.  On \(\mathrm dz\),
the first three terms telescope to
\(q^{r+1}z(T_q)-z(T_q^q)\); the inertia sum is
\(qz(T_q)-z(T_q^q)\).  Division by \(q\), a unit at
\(\ell\), proves the literal identity
\[
 \psi_q^2(\mathrm dz)=(q^r-1)z(T_q),\qquad
 \psi_q^2(R)=-(q^r-1)z_q(T_q).
\tag{A.204.11}
\]
The discarded inertia acts trivially and maps trivially
to the pro-\(\ell\) geometric group, so these cochains
factor through the actual tame quotient, as in A.201.

Put \(R_\ell^0=\{(w_q)_{q\ne\ell}:w_q\in\Z_\ell,
w_q\to0\}\).  The localization sequence (A.65.8) is
exact at its middle term and has final target
\[
          (R_\ell^0)^{b_{r+1}}/
             (q^r-1)(R_\ell^0)^{b_{r+1}}.
\tag{A.204.12}
\]
Its denominator is the literal image.  Equations
(A.204.10)--(A.204.11) place the residue of the original
global defect in that image, with explicit primitive
\((-z_q(T_q))_q\).  Exactness proves (A.204.1), including
the continuous inverse-limit contribution.
The already proved annihilator (A.65.7), with weight
\(r+1\), gives (A.204.2).  Its precise inputs are
\cite[VI, Theorems 8.2 and 9.11]{Weibel2013}, with
the finite positive even \(K\)-group and the retained
real two-primary term; no uncomputed \(K\)-group is
declared zero here.

\emph{Vanishing in weight four.}
Set \(r=3\).  The exact computation
(A.61.18)--(A.61.20), using
\cite[VI, Theorems 8.2, 9.11 (case \(n=6\)),
and Examples 10.3]{Weibel2013}, gives
\[
\begin{aligned}
 H^2_{\rm cts}(G_{\Q,\{\ell,\infty\}},\Z_\ell(4))&=0
                                      &&(\ell\text{ odd}),\\
 H^2_{\rm cts}(G_{\Q,\{2,\infty\}},\Z_2(4))
       &\hookrightarrow H^2(G_\R,\Z_2(4)).
\end{aligned}
\tag{A.204.13}
\]
The kernel calculation is exactly \(K_6(\Z)=0\),
including its two-primary part.  At two the real target
\(\Z/2\) is retained.  By the assumed actual real full
point or cusp lift, the global obstruction restricts
to zero there.  Any unramified preimage supplied by
(A.204.1) therefore has zero real restriction, because
restriction commutes with inflation to \(G_\Q\).
The injection in (A.204.13) makes that preimage zero.
All three obstruction coordinates consequently vanish
at every coefficient prime.

Choose continuous central primitives of \(-R\) in these
actual coefficient complexes.  Their product is continuous,
since continuous cochains with a product of compact Tate
modules are exactly the product of the coefficient cochains.
The geometric nilpotent quotient decomposes into its
characteristic pro-\(\ell\) factors as in A.201.
Correcting the highest Hall coordinates by that product
therefore gives an actual global section of \(E_4\)
projecting to the prescribed representative of \(s\).
Lifting geometric conjugators proves (A.204.3) on section
classes.  This construction solves the global obstruction;
it makes no assertion about the local correction classes
of the resulting lift.

The finite certificate
\path{code/selmer_weight_four_inertia_certificate.py}
checks the literal fourth-layer Hall extraction and all
three cusp conjugations.  The all-weight uniform limit
is the compactness proof (A.204.8)--(A.204.10), and the
global lifting assertion uses the strict residue and
unramified maps above.
\end{proof}
